\section{Conductors and local models of divisor collisions}
\label{sec:conductor-v152}

We retain the notation of Section~\ref{sec:canonical-boundary-v151}.
All completed local rings below are completed at closed complex points.
For a reduced local ring $A$ with finite normalization $B$, its
\emph{conductor} is $\mathfrak c=\Ann_A(B/A)$, regarded also as an
ideal of $B$.  We use the intrinsic singular subscheme defined by
$\Fitt_{\dim A}\Omega_{A/\C}$; on completions, differentials mean
continuous differentials.  The conductor and this Fitting ideal need
not agree.  Our first result determines both on a multivariate
stratum.  The second identifies the way two branches merge when a
common binary root becomes double.

\subsection{Two coprime divisor choices}
\begin{theorem}[The split-divisor stratum]\label{thm:split-conductor-v152}
Suppose $h>g$, $2\leq r\leq s_{h-g}$, and
\[
 K=abL,\qquad [a],[b]\in\Pj(S_g),\qquad
 L\in\Gr(r,S_{h-g}),\qquad \gcd(L)=1.
\]
Assume that $a,b$ are coprime and that the only degree-$g$ divisors of
$ab$, up to scalar, are $a$ and $b$.  Put
\begin{align*}
 t&=2(s_g-1)+r(s_{h-g}-r),\\
 c&=r(s_h-s_{h-g})-(s_g-1),\qquad d_g=t+c.
\end{align*}
Then $c\geq1$, and there are formal coordinates in which
\begin{equation}\label{eq:split-ring-v152}
 \widehat{\OO}_{Y_g,K}=
 \C[[z_1,\ldots,z_t,x_1,\ldots,x_c,y_1,\ldots,y_c]]/
                       (x_i y_j:1\leq i,j\leq c).
\end{equation}
Its normalization is
$\C[[z,x]]\times\C[[z,y]]$.  The conductor is
\begin{equation}\label{eq:split-conductor-v152}
 \mathfrak c=(x_1,\ldots,x_c,y_1,\ldots,y_c),\qquad
 \widehat{\OO}_{Y_g,K}/\mathfrak c=\C[[z]],
\end{equation}
and the intrinsic singular subscheme has ideal $\mathfrak c^{\,c}$.
The two smooth branches meet scheme-theoretically in a smooth
$t$-dimensional subscheme.  The local ring has multiplicity two and
depth $t+1$; it is Cohen--Macaulay exactly when $c=1$.
\end{theorem}

\begin{proof}
We first determine the intersection as a scheme, not by comparing
closed points or tangent dimensions.  The normalization has precisely
two points above $K$, namely $(a,bL)$ and $(b,aL)$.  At either point,
the overlap in Lemma~\ref{lem:fibre-equations-v151} is zero.  The
map from the completed local ring of the ambient Grassmannian to
each completed source ring is consequently surjective: its cotangent
map is surjective, and completeness and Nakayama's lemma lift
successive homogeneous terms.  Thus each source completion is a
closed smooth formal branch of the image.

Here is the required assertion over infinitesimal bases.  Let $A_0$
be a local Artin $\C$-algebra.  A point on both formal branches carries
unique marked lifts $a_{A_0},b_{A_0}$ near $a,b$.  If homogeneous
polynomials have coprime reductions, those reductions form a regular
sequence in the polynomial ring over $\C$.  Its graded Koszul
complex is exact.  In each degree the polynomial modules are finite
free; lifting a splitting of the complex, or applying the local
criterion for flatness successively, shows that its lift over $A_0$
is exact and that the quotient is $A_0$-flat.  In particular
\begin{equation}\label{eq:coprime-artin-v152}
 (a_{A_0})\cap(b_{A_0})=(a_{A_0}b_{A_0}).
\end{equation}
This also follows by noting that $b_{A_0}$ is a nonzerodivisor
modulo $a_{A_0}$ and solving a relation between them.  Applying
\eqref{eq:coprime-artin-v152} to the columns of the relation
subbundle gives, uniquely,
\[
                  K_{A_0}=a_{A_0}b_{A_0}L_{A_0}.
\]
The residual module is a subbundle: multiplication by
$a_{A_0}b_{A_0}$ is a split injection in degree $D$, and the
quotient of the ambient free module by $K_{A_0}$ is free.  The
kernel of the induced surjection onto the complementary quotient
is therefore free.  This proves that the formal intersection is
represented by the product of the two divisor charts and the
Grassmannian chart for $L$, with no nilpotents omitted.

The differential of this product parametrization is injective.
Indeed, an infinitesimally constant product gives
\[
       (\dot a b+a\dot b)l+ab\dot l\in abL\quad(l\in L).
\]
Taking valuations at the factors of $a$, and using
$\gcd(a,b)=\gcd(L)=1$, forces $a\mid\dot a$; similarly
$b\mid\dot b$.  These are the two projective scalar directions,
and then $\dot L=0$.  The intersection is thus a smooth formal
closed subscheme of dimension $t$.

Let $R_0$ be the ambient completed local ring and $I_a,I_b$ the
ideals of these branches.  Finiteness of the normalization, exactness
of completion on finite modules, and the definition of the
scheme-theoretic image give
\[
 \widehat{\OO}_{Y_g,K}=R_0/(I_a\cap I_b)
   = (R_0/I_a)\mathbin{\times}_{R_0/(I_a+I_b)}(R_0/I_b).
\]
This is the elementary conductor-square construction; see also
\cite[Lemma~1.3 and \S1.3.2]{Ferrand2003}.  All three quotient rings
are regular.  Choosing compatible regular parameters for the two
surjections gives \eqref{eq:split-ring-v152}.  The branch dimension
is $d_g=s_g-1+r(s_h-r)$, so its codimension along the intersection
is the asserted $c$.  It is positive: if $c=0$, both branch ideals
would coincide locally with the intersection ideal, contradicting
the two distinct normalization branches of a reduced local ring.

In the fibre-product description, an element belongs to the
conductor exactly when both its restrictions to the intersection
vanish.  This gives \eqref{eq:split-conductor-v152}.  The Jacobian
of the relations $x_i y_j$ has nonzero entries only among the $x_i$
and $y_j$.  The Fitting ideal in question is its ideal of $c$-minors.
Every such minor belongs to $\mathfrak c^{\,c}$.  Conversely,
choosing the $c$ columns indexed by the $y_j$ and one row $(i_j,j)$
for each $j$ gives every monomial $x_{i_1}\cdots x_{i_c}$.
The symmetric choice gives every degree-$c$ monomial in the $y$'s.
Mixed monomials vanish in the ring.  Thus the ideal is exactly
$\mathfrak c^{\,c}$, including its nonreduced structure.

Finally use
\[
 0\longrightarrow \widehat{\OO}_{Y_g,K}
 \longrightarrow\C[[z,x]]\oplus\C[[z,y]]
 \longrightarrow\C[[z]]\longrightarrow0,
\]
where the last arrow is the difference of restrictions.  The depth
lemma gives depth $t+1$; for $c=1$ this equals the dimension.
For $c>1$ it is strictly smaller.  Additivity of leading Hilbert
coefficients gives multiplicity two, since the last module has
smaller dimension than the two regular branches.
\end{proof}

\begin{remark}\label{rem:split-scope-v152}
The hypotheses are nonempty in every number of variables: for $e=2$
one may take $g=1$, and for $e\geq3$ one may take two distinct
irreducible forms of degree $g$.  The theorem treats an explicitly
specified stratum, not all pairs of overlapping degree-$g$ divisors.
In particular, it does not substitute a binary factorization rule
for multivariate divisibility.  Its high-codimension branch
intersections need not be normal-crossing divisors.
\end{remark}

\subsection{A double common root}
For $m\geq1$, put $R_m=\C[[u_1,\ldots,u_m,\Delta]]$ and
$I=(u_1,\ldots,u_m)\subset R_m$.  Define
\begin{equation}\label{eq:pinch-algebra-v152}
 A_m=R_m\oplus I t\ \subset\ B_m=\C[[u_1,\ldots,u_m,t]],
                    \qquad \Delta=t^2.
\end{equation}
The right side of \eqref{eq:pinch-algebra-v152} is a subring,
not a square-zero extension: $(It)^2=\Delta I^2$.  For $m=1$
it is the ordinary pinch-point ring.

\begin{theorem}[Collision of two binary divisor choices]
\label{thm:binary-pinch-v152}
Let $e=2$, $g=1$, $D=h+1$, and $2\leq r\leq D-1$.
At a point $K_0=\ell^2 L$ with
$L\in\Gr(r,S_{D-2})$ and $\gcd(L)=1$, set
\[
             m=r-1,\qquad q=r(D-1-r)+1.
\]
Then
\begin{equation}\label{eq:binary-completion-v152}
              \widehat{\OO}_{Y_1,K_0}\simeq A_m[[z_1,\ldots,z_q]].
\end{equation}
Writing $w_i=u_i t$, a complete presentation of $A_m$ is
\begin{equation}\label{eq:pinch-equations-v152}
 \C[[u_1,\ldots,u_m,\Delta,w_1,\ldots,w_m]]/
 \bigl(u_iw_j-u_jw_i,\ w_iw_j-\Delta u_i u_j\bigr),
\end{equation}
where the first relations have $i<j$ and the second have $i\leq j$.
Its normalization is $B_m$, and its conductor is
\[
 \mathfrak c=(u_1,\ldots,u_m,w_1,\ldots,w_m),\qquad
 \mathfrak c B_m=I B_m.
\]
The induced map on conductor quotients is
\begin{equation}\label{eq:conductor-cover-v152}
              \C[[\Delta]]\longrightarrow\C[[t]],\qquad\Delta\longmapsto t^2.
\end{equation}
Thus the two reduced divisor lifts merge through a ramified double
cover of the conductor.  The completed local ring has dimension
$q+m+1$ and depth $q+2$.  It is Cohen--Macaulay if and only if $r=2$.
For $r=2$, its intrinsic singular ideal, with the smooth variables
suppressed, is $(w,\Delta u,u^2)$, whereas its conductor is $(u,w)$.
\end{theorem}

\begin{proof}
Choose an affine coordinate $x$ with the root of $\ell$ at zero.
Some member $F_0\in K_0$ has multiplicity exactly two there, since
$\gcd(L)=1$.  Extend it to a frame $F_0,\ldots,F_{r-1}$.
In the completed local ring of the ambient Grassmannian the formal
Weierstrass theorem writes the universal $F_0$ as a unit times
\[
                         P(x)=x^2+a x+b.
\]
Divide each $F_i$, $i>0$, by this monic polynomial.  Its remainder
is $u_i x+v_i$.  The $2r$ functions
$a,b,u_1,v_1,\ldots,u_m,v_m$ are part of a regular parameter
system.  To see this rather than assume versality, note that
$K_0$ has zero two-jet at the root, while
$S_D/K_0\to\C[x]/(x^2)$ is surjective.  The Grassmannian tangent
space $\Hom(K_0,S_D/K_0)$ therefore allows independent variations
of both jets in every frame column.  Passing from these jets to
the displayed remainders is an invertible triangular operation,
with leading coefficient $F_0/x^2(0)\ne0$.  Their differentials
are consequently independent.

Set
\[
       t=x+a/2,\quad \Delta=a^2/4-b,\quad
       w_i=a u_i/2-v_i.
\]
The incidence of a common root near zero is now exactly
\[
                       \Delta=t^2,\qquad w_i=u_i t.
\]
There are no other normalization points above $K_0$, because its
exact common divisor is $\ell^2$.  Consequently the image of this
completed incidence is the completed image defining $Y_1$.
The ambient dimension is $N_0=r(D+1-r)$.  Of the $2r$ parameters
above, $a$ is free, as are the other $N_0-2r$ parameters.  This
accounts for $q=N_0-2r+1$ smooth variables in
\eqref{eq:binary-completion-v152}.

It remains to justify elimination scheme-theoretically.  The
relations in \eqref{eq:pinch-equations-v152} reduce any series to
$p+\sum_i w_i p_i$ with $p,p_i\in R_m$.  If its image in $B_m$
is zero, the decomposition $B_m=R_m\oplus R_m t$ gives
$p=0$ and $\sum_i u_i p_i=0$.  The first syzygies of the regular
sequence $u_1,\ldots,u_m$ are generated by the Koszul syzygies.
They are precisely the relations $u_iw_j-u_jw_i$.  This proves
the presentation without taking a radical.  The same argument
works first over polynomial rings and then by exact completion.

The extension $A_m\subset B_m$ is finite, since $t^2=\Delta$,
and birational, since $t=w_1/u_1$ in the fraction field.  The
regular ring $B_m$ is thus the normalization.  For
$p+v t\in R_m\oplus I t$, multiplication by $t$ stays in
$A_m$ if and only if $p\in I$.  As $1,t$ generate $B_m$,
the conductor is $I\oplus I t$, proving the quotient map
\eqref{eq:conductor-cover-v152}.  Equivalently,
\begin{equation}\label{eq:pinch-square-v152}
 A_m=B_m\mathbin{\times}_{\C[[t]]}\C[[\Delta]],
\end{equation}
where the first map sets $u=0$ and the second sends $\Delta$ to
$t^2$.  This is a specific conductor square, not an assumption
that the normalization descends through its reduced fibres.

As an $R_m$-module, $A_m=R_m\oplus I$.  For $m=1$, the ideal
$I$ is free.  For $m>1$, the exact sequence
$0\to I\to R_m\to\C[[\Delta]]\to0$ gives
$\operatorname{depth}_{R_m}I=2$.  Hence $A_m$ has depth two
for every $m\geq1$.  The maximal ideal of $R_m$ generates an
ideal with maximal radical in $A_m$, so this is also its local
ring depth.  Adding the $q$ smooth parameters gives the assertion.
When $m=1$, the equation is $w^2-\Delta u^2=0$; its partial
derivatives generate $(w,\Delta u,u^2)$ in characteristic zero.
\end{proof}

\begin{corollary}[The split and ramified parts of the same model]
\label{cor:collision-strata-v152}
On the binary locus where the common divisor has degree exactly two,
there are two strata.  Over two distinct roots the formal local ring
has the split form \eqref{eq:split-ring-v152}, with
$t=q+1$ and $c=r-1$.  Over a double root it has
\eqref{eq:binary-completion-v152}.  The support of the singular
locus in the latter model is the conductor, and the normalization
fibre at a collision is $\Spec\C[t]/(t^2)$.
For arbitrary $m$, its singular subscheme is given without
radicalization by the $m$-minors of the Jacobian of the displayed
$m^2$ relations in \eqref{eq:pinch-equations-v152}.
\end{corollary}
\begin{proof}
The assertion at distinct roots is Theorem~\ref{thm:split-conductor-v152}.
In the collision model the normalization is an isomorphism wherever
some $u_i$ is invertible.  On $u=w=0$, a nonzero value of $\Delta$
gives two distinct points after an \'etale square-root extension,
while $\Delta=0$ gives the stated length-two fibre.  The exact
normality criterion proves the assertion about singular support.
There are $2m+1$ nonsmooth coordinates and dimension $m+1$;
the cotangent presentation therefore gives the $m$-minors as
claimed.  This formula specifies the whole singular scheme, even
where it is thicker than its conductor support.
\end{proof}

\begin{remark}[Descent to intrinsic first relations]
\label{rem:conductor-stack-v152}
On the maximal-lower-Hilbert-function stratum of
Corollary~\ref{cor:algebra-boundary-v151}, these statements compute
the completed smooth-atlas singularities of the intrinsic
first-relation algebra stack after tangent-identity rigidification.
The conductor is equivariant, and formation of its finite-module
annihilator and of the Fitting ideals commutes with flat base change.
Thus the calculations descend through the frame atlas.  This does
not identify these rings with completed local rings of a coarse
orbit space, where stabilizers can introduce further singularities.
\end{remark}
